\section*{Chapitre \ref{ch:g8:fronts-and-weather-maps} --- Masses d'air, fronts et cartes météorologiques}

\begin{solution}{exo:g8:fronts-and-weather-maps:1}
(a) Maritime tropicale~: chaude et humide. (b) Continentale polaire~:
froide et sèche. (c) Continentale tropicale~: très chaude et sèche.
(d) Maritime polaire~: froide et humide.
\end{solution}

\begin{solution}{exo:g8:fronts-and-weather-maps:2}
Ligne bleue à triangles~: un \omterm{def:g8:fronts-and-weather-maps:front}{front froid}. Ligne rouge à demi-cercles~: un
\omterm{def:g8:fronts-and-weather-maps:front}{front chaud}. Ligne qui joint des lieux de même pression~: une \omterm{def:g8:fronts-and-weather-maps:isobar}{isobare}.
A~: le centre d'un \omterm{def:g8:fronts-and-weather-maps:depression}{anticyclone} (hautes pressions). D~: le centre d'une
\omterm{def:g8:fronts-and-weather-maps:depression}{dépression} (basses pressions).
\end{solution}

\begin{solution}{exo:g8:fronts-and-weather-maps:3}
(a) Vrai. (b) Faux~: ils arrivent avant lui, loin en avant, haut dans le
ciel. (c) Faux~: des \omterm{def:g8:fronts-and-weather-maps:isobar}{isobares} serrées signifient un vent fort. (d) Vrai.
\end{solution}

\begin{solution}{exo:g8:fronts-and-weather-maps:4}
(a) \qty{980}{hPa}. (b) \qty{1034}{hPa}. (c) \qty{4}{hPa}. (d) $1034 -
980 = \qty{54}{hPa}$.
\end{solution}

\begin{solution}{exo:g8:fronts-and-weather-maps:5}
De fines traînées de \omterm{def:g4:clouds-rain-snow:cloud-families}{cirrus}~; un voile gris d'altostratus~; les \omterm{def:g4:clouds-rain-snow:cloud}{nuages} qui
descendent et s'épaississent~; une \omterm{def:g4:clouds-rain-snow:precipitation}{pluie} continue~; un air plus chaud et
plus humide.
\end{solution}

\begin{solution}{exo:g8:fronts-and-weather-maps:6}
À un \omterm{def:g8:fronts-and-weather-maps:front}{front froid}, l'air froid et dense se glisse sous l'air chaud comme un
\emph{coin}~: l'air chaud est soulevé \emph{brusquement}, sur une pente raide,
monte vite et forme de hauts \omterm{def:g4:clouds-rain-snow:cloud}{nuages}. À un \omterm{def:g8:fronts-and-weather-maps:front}{front chaud}, l'air chaud monte
lentement par-dessus l'air froid, sur une \emph{pente douce}~: il s'élève
peu à peu sur une longue distance, et ses \omterm{def:g4:clouds-rain-snow:cloud}{nuages} s'étalent en couches.
\end{solution}

\begin{solution}{exo:g8:fronts-and-weather-maps:7}
(a) Près de D~: les \omterm{def:g8:fronts-and-weather-maps:isobar}{isobares} y sont plus serrées, et les flèches plus
longues. (b) De l'est (les flèches pointent vers l'ouest, un peu vers le
sud-ouest). (c) De l'est (les flèches pointent vers l'ouest). (d) Dans le
sens inverse des aiguilles d'une montre.
\end{solution}

\begin{solution}{exo:g8:fronts-and-weather-maps:8}
(a) $1016 - 1003 = \qty{13}{hPa}$. (b) Une \omterm{def:g8:fronts-and-weather-maps:depression}{dépression}~: la pression ne
cesse de baisser. (c) Des \omterm{def:g4:clouds-rain-snow:cloud}{nuages}, du vent et de la \omterm{def:g4:clouds-rain-snow:precipitation}{pluie}, probablement
avec un \omterm{def:g8:fronts-and-weather-maps:front}{front}.
\end{solution}

\begin{solution}{exo:g8:fronts-and-weather-maps:9}
(a) cT, continentale tropicale. (b) Très chaude et sèche. (c) Au-dessus de
la \omterm{def:g2:water-on-earth:ocean}{mer}, elle se charge de vapeur d'eau~: elle devient plus humide, et l'air
proche de l'eau est un peu refroidi par la \omterm{def:g2:water-on-earth:ocean}{mer}.
\end{solution}

\begin{solution}{exo:g8:fronts-and-weather-maps:10}
(a) $4 - 1 = 3$~jours. (b) $3 \div 3 = 1$~jour tous les dix ans. (c) Il
faut encore $5 - 1 = 4$~jours~: environ 40~ans.
\end{solution}

\begin{solution}{exo:g8:fronts-and-weather-maps:11}
(a) Environ \qty{1350}{km} (quatre fois et demie l'\omterm{def:g2:drawing-to-map:scale}{échelle} graphique). (b) \qty{4}{hPa}
sur \qty{80}{km}~: $4 \div 80 \times 100 = \qty{5}{hPa}$ pour
\qty{100}{km}. (c) $4 \div 140 \times 100 \approx \qty{2.9}{hPa}$ pour
\qty{100}{km}. (d) Près de D, où la pression change plus vite~; oui, les
flèches sont plus longues près de D.
\end{solution}

\begin{solution}{exo:g8:fronts-and-weather-maps:12}
(a) Le \omterm{def:g8:fronts-and-weather-maps:front}{front froid}~: $300 \div 50 = 6$~heures. Le \omterm{def:g8:fronts-and-weather-maps:front}{front chaud} avance à
\qty{25}{km} par heure~: $200 \div 25 = 8$~heures. Le \omterm{def:g8:fronts-and-weather-maps:front}{front froid} arrive
2~heures plus tôt. (b) Première ville~: des \omterm{def:g4:clouds-rain-snow:cloud}{nuages} en tours, de fortes
averses, des rafales, puis un air plus frais et un ciel qui se dégage.
Seconde ville~: des \omterm{def:g4:clouds-rain-snow:cloud}{nuages} en couches qui descendent et une longue \omterm{def:g4:clouds-rain-snow:precipitation}{pluie}
continue, puis un air plus chaud et plus humide. (c) La seconde ville,
atteinte par l'air chaud.
\end{solution}

\begin{solution}{pb:g8:fronts-and-weather-maps:1}
\textbf{1.} Dans l'air chaud, entre les deux \omterm{def:g8:fronts-and-weather-maps:front}{fronts}. Le \omterm{def:g8:fronts-and-weather-maps:front}{front chaud} est
déjà passé.

\textbf{2.} Du sud-ouest.

\textbf{3.} Vers l'est~: les triangles d'un \omterm{def:g8:fronts-and-weather-maps:front}{front froid} sont placés du côté
vers lequel il avance, et ils pointent vers l'est, vers la ville.

\textbf{4.} Des \omterm{def:g4:clouds-rain-snow:cloud}{nuages} en tours, de fortes averses, peut-être du tonnerre
et de la \omterm{def:g4:clouds-rain-snow:precipitation}{grêle}, des rafales de vent~; puis un air plus frais, un vent
d'ouest ou de nord-ouest et un ciel qui se dégage.

\textbf{5.} À \qty{180}{km} à 6~h et à \qty{300}{km} à 9~h.

\textbf{6.} $300 - 180 = \qty{120}{km}$.

\textbf{7.} $120 \div 3 = \qty{40}{km}$ par heure.

\textbf{8.} La ville est à \qty{480}{km}~: $480 - 300 = \qty{180}{km}$.

\textbf{9.} $180 \div 40 = 4.5$~heures, soit 4~heures et 30~minutes.

\textbf{10.} 9~h + 4~h~30~min = \textbf{13~h~30}.

\textbf{11.} Non~: les averses devraient se terminer vers 14~h~30, avant
le coup d'envoi.

\textbf{12.} $180 \div 30 = 6$~heures~: à 15~h, juste au coup d'envoi.
Prévision~: «~Le \omterm{def:g8:fronts-and-weather-maps:front}{front froid} devrait atteindre la ville vers
\textbf{13~h~30}, avec une heure de fortes averses et de rafales, puis
des éclaircies~: le match pourra très probablement se jouer. Si le \omterm{def:g8:fronts-and-weather-maps:front}{front}
ralentit, il pourrait arriver à 15~h et retarder le match.~»
\end{solution}
